% =============================================================================
% §1.8 Le calcul programmable et la révolution électronique (1945–1975)
% =============================================================================
% Résultats dimensionnels visés :
%   0-TER  — le computing comme problème énergétique visible
%   D3     — séparation hardware/software, naissance du capital intangible
%   D2     — consent decree comme anti-monopole (négatif du télégraphe)
%   D5/D8  — État comme créateur de marché
%   D10    — Moore's Law comme dispositif de coordination sectorielle
% Sources empiriques principales :
%   Nordhaus 2007, Koomey et al. 2011, Dennard et al. 1974,
%   Flamm 1987/1988, Moore 1965, Mody 2017, Merges & Nelson 1990
% =============================================================================

\section{Le calcul programmable et la révolution électronique (1945--1975)}\label{sec:computing}

Les deux cas précédents ont mis en lumière un contraste de fonction et de trajectoire. Le télégraphe transmet de l'information entre des lieux distants, à un coût énergétique si faible par rapport au transport physique qu'il en est invisible ; la mécanisation du calcul traite de l'information sur place, à un coût en travail humain que la machine redistribue sans le supprimer. Dans les deux cas, l'énergie consommée par le traitement de l'information reste un poste négligeable : 160 watts pour un bureau télégraphique contre 370 kilowatts pour la locomotive qui acheminait les mêmes dépêches ; neuf livres sterling d'électricité sur un budget annuel de six cent sept livres au Nautical Almanac Office mécanisé.

Avant que le calcul électronique n'existe matériellement, son concept est posé. \textcite{turing_computable_1936} formule, dans \emph{On Computable Numbers} paru aux \emph{Proceedings of the London Mathematical Society} en 1937, le concept de la machine universelle~--- une machine unique capable d'exécuter tout calcul reproductible pourvu qu'on lui fournisse la description du calcul à effectuer\footnote{Le mot «~ordinateur~» étant postérieur à cette formulation et propre à la langue française (1955, sur proposition de Jacques Perret pour IBM France), Turing parle de \emph{computing machine}, dont sa version universelle est la matrice conceptuelle. La projection rétrospective du mot ordinateur sur le texte de 1937 est un anachronisme léger, sans conséquence pour l'argument.}. Le résultat est purement mathématique~: aucun dispositif matériel n'en est tiré à ce moment. Neuf ans séparent cette formulation théorique de l'ENIAC, et quinze ans la séparent du premier ordinateur commercial~; l'écart entre la formulation conceptuelle et la disponibilité industrielle, caractéristique des techniques de l'information qui se reproduira à plusieurs reprises dans ce chapitre, ne tient pas à une difficulté technique surmontée par les ingénieurs, mais à l'absence préalable d'une demande qui aurait justifié la mobilisation des ressources nécessaires à la matérialisation du concept. La généalogie qui mène d'ENIAC à EDVAC puis à l'\emph{IAS Machine} peut donc se lire à deux niveaux indépendants~: une généalogie matérielle dont les jalons sont des artefacts physiques, et une généalogie conceptuelle dont les jalons sont des publications théoriques. Le présent chapitre suit les deux.

\dimmark{0-TER}{*}%
Le cas qui s'ouvre ici rompt avec cette invisibilité. En 1946, l'Electronic Numerical Integrator and Computer consomme à lui seul 150 kilowatts, soit la puissance d'une petite centrale, pour exécuter cinq mille additions par seconde. Pour la première fois dans l'histoire du traitement de l'information, la consommation d'énergie de la machine est du même ordre de grandeur que celle des activités qu'elle sert. Ce n'est pas un accident de jeunesse technologique ; c'est le début d'une course de quatre-vingts ans entre la puissance de calcul et la consommation électrique qui la rend possible, course dont les termes se modifient à chaque génération technologique sans que le couplage structurel se desserre.

Cette visibilité de l'énergie est cependant d'un type particulier. Elle est visible comme \emph{contrainte d'ingénierie}~: l'ENIAC nécessite une alimentation dédiée, un refroidissement actif par ventilateurs, une tension régulée pour éviter la corruption des calculs par les fluctuations du réseau\footnote{L'ENIAC opère sur dix-huit mille tubes à vide dont chacun peut tomber en panne --- en moyenne un tube toutes les deux heures en fonctionnement continu. La stabilité de l'alimentation électrique conditionne directement la fiabilité du calcul. Ce couplage entre qualité de l'énergie et intégrité de l'information est une propriété structurelle des premiers ordinateurs électroniques, non un détail opérationnel.}. Elle est en revanche \emph{absente des comptabilités nationales et des modèles économiques de la productivité}~: la \emph{Survey of Current Business} n'isole pas de poste «~énergie du traitement de l'information~»; les travaux de \textcite{flamm_creating_1988} sur le financement de l'informatique ne mentionnent pas le coût électrique d'exploitation. La visibilité technique de l'énergie coexiste avec son invisibilité économique et statistique. C'est cette dissymétrie~--- et non l'ignorance brute d'une contrainte~--- que la thèse cherche à documenter.

\dimmark{0-TER}{+}%
Cette dissymétrie a une justification physique que les travaux de Landauer ont établie dès 1961. \textcite{landauer_irreversibility_1961} montre que l'effacement d'un bit dans une mémoire physique dissipe nécessairement une quantité minimale d'énergie de l'ordre de $kT\ln 2$ joules, où $k$ est la constante de Boltzmann et $T$ la température du système~--- environ $3 \times 10^{-21}$~joules à température ambiante\footnote{La borne de Landauer a été confirmée expérimentalement par \textcite{berut_experimental_2012} et par \textcite{jun_high_2014}. \textcite{bennett_thermodynamics_1982} a montré que le calcul réversible (sans effacement d'information) échappe en théorie à cette borne~; en pratique, tout ordinateur réel efface de l'information et la borne s'applique. Le ratio entre la consommation énergétique effective des processeurs modernes et la borne de Landauer est de l'ordre de $10^{6}$~--- autant de marge d'amélioration technologique encore disponible, et autant de preuve que la consommation actuelle tient à l'architecture des composants, pas aux lois de la physique.}. Ce résultat est fondamental pour la thèse~: il établit que le coût énergétique du calcul n'est pas un accident technologique corrigible par l'ingénierie~--- il est irréductible au niveau du bit. La convergence des trois fonctions sur un substrat numérique après 1983 ne produit pas seulement un coût énergétique économiquement visible~; elle produit un couplage physiquement nécessaire entre traitement de l'information et dissipation d'énergie.

Ce cas couvre trente ans, de l'ENIAC (1946)\footnote{La paternité du \emph{stored-program concept} --- le principe selon lequel le programme est stocké dans la même mémoire que les données --- est disputée dans l'historiographie. Konrad Zuse présente le Z3 en 1941, premier ordinateur programmable fonctionnel à relais électromécaniques, dont les autorités nazies sous-estimeront l'importance. Eckert et Mauchly (ENIAC/EDVAC) revendiquent la priorité sur Von Neumann, dont le \emph{First Draft of a Report on the EDVAC} (1945) formalise l'architecture à programme enregistré de façon si complète que son nom reste attaché au concept. Ce qui importe analytiquement pour ce cas n'est pas la priorité d'invention mais la conséquence architecturale~: l'espace mémoire programme et l'espace mémoire données sont unifiés dans la même structure. C'est cette fusion~--- et non l'existence d'un ordinateur programmable~--- qui marque le pivot de la convergence STORE/COMPUTE.} au microprocesseur Intel 4004 (1971) et au papier fondateur du \emph{MOSFET scaling} de Dennard (1974). Il traverse au moins quatre artefacts majeurs : l'ordinateur programmable à tubes, le mainframe commercial à transistors, le circuit intégré, le microprocesseur. L'unité du cas ne tient pas à un objet technique unique, comme le télégraphe de Morse ou la tabulatrice de Hollerith, mais à une dynamique commune : chaque bond de performance computationnelle ouvre un nouveau front de consommation énergétique, tandis que chaque gain d'efficacité par composant est réinvesti en densité de calcul. \textcite{nordhaus_two_2007} mesure la chute du coût par calcul sur deux siècles ; \textcite{koomey_implications_2011} mesure le doublement de l'efficacité énergétique par calcul tous les dix-huit mois entre 1946 et 2000. Personne ne croise les deux séries. C'est dans cet écart que la thèse se loge.


% -----------------------------------------------------------------------------
\subsection{Du tube au transistor~: la transition matérielle (1947--1965)}\label{subsec:transistor}
% -----------------------------------------------------------------------------

\dimmark{D10}{+}%
Le 23~décembre 1947, aux Bell Telephone Laboratories, John Bardeen, Walter Brattain et William Shockley réalisent le premier transistor à pointe de contact. Le dispositif, encore fragile et coûteux, remplace le tube à vide dans sa fonction d'amplification et de commutation~: plus petit (quelques millimètres contre plusieurs centimètres), plus fiable (pas de filament à chauffer), et surtout beaucoup moins gourmand en énergie (milliwatts contre watts). Le brevet est déposé en 1948~; le prix Nobel couronne les trois inventeurs en 1956. La trajectoire qui va du premier transistor au circuit intégré couvre moins de quinze ans, une compression temporelle sans équivalent dans les cas précédents.

\dimmark{D5/D8}{+}%
Cette compression n'est pas le fait du marché civil. \textcite{flamm_creating_1988} documente que la demande militaire absorbe la quasi-totalité de la production de transistors dans les premières années. Le programme de miniaturisation des missiles balistiques intercontinentaux (Minuteman) exige des composants légers et fiables~; la Navy et l'Air Force achètent 100\,\% de la production américaine de circuits intégrés en 1962. La courbe d'apprentissage qui en résulte fait chuter le prix du transistor de 45~dollars en 1954 à 2~dollars en 1960~--- une division par vingt en six ans. Le marché civil hérite d'une technologie dont le coût de développement a été socialisé par le budget de la défense.

\dimmark{D10}{+}%
Le circuit intégré (1958--1965) franchit une étape supplémentaire. Jack Kilby chez Texas Instruments et Robert Noyce chez Fairchild Semiconductor déposent des brevets indépendants la même année (1958-1959) pour l'intégration de plusieurs transistors sur un même substrat de silicium. La querelle de priorité qui s'ensuit~--- tranchée par un accord de licences croisées en 1966~--- importe moins que la conséquence architecturale~: le circuit intégré permet de graver des milliers, puis des millions de transistors sur une puce unique, réduisant les connexions externes et les points de défaillance. Noyce fonde Intel en 1968 avec Gordon Moore~; le premier microprocesseur commercial (Intel 4004) sort en 1971. La trajectoire de la loi de Moore (§\,\ref{subsec:moore_law}) commence ici.

\dimmark{0-TER}{+}%
Le profil énergétique de cette transition est paradoxal. L'efficacité par opération s'améliore de façon spectaculaire~: le ratio joules/opération chute de six ordres de grandeur entre l'ENIAC (1946) et l'IBM~360 (1964)\footnote{L'ENIAC consommait environ 150\,kW pour 5\,000 additions par seconde, soit 30~joules par opération. L'IBM~360/50 (1964) consommait environ 20\,kW pour 500\,000 opérations par seconde, soit 0,04~joule par opération. Le gain est de l'ordre de $10^{3}$. Entre l'IBM~360 et l'Intel~4004 (1971), le gain est encore de trois ordres de grandeur. Sources~: \textcite{nordhaus_two_2007}, \textcite{koomey_implications_2011}.}. Mais la puissance totale installée augmente~: les mainframes des années 1960 consomment 20 à 50~kW chacun, et leur nombre croît. Le paradoxe est déjà en place~: l'efficacité par unité s'améliore, la consommation agrégée croît parce que le volume de calcul croît plus vite que l'efficacité. C'est le mécanisme que \textcite{jevons_coal_1865} avait identifié pour le charbon et que la littérature contemporaine nomme \emph{rebound effect}~; il s'appliquera au calcul pendant les cinq décennies suivantes.


% -----------------------------------------------------------------------------
\subsection{La demande de calcul comme problème militaire}\label{subsec:demand_pull_computing}
% -----------------------------------------------------------------------------

\dimmark{D5}{+}%

La demande qui fait naître le calcul électronique n'est ni commerciale ni statistique. Elle est militaire. En 1942, le Ballistic Research Laboratory d'Aberdeen, dans le Maryland, doit produire les tables de tir de l'artillerie américaine : chaque table exige le calcul de plusieurs centaines de trajectoires, chaque trajectoire requiert la résolution numérique d'équations différentielles non linéaires tenant compte de la gravité, de la résistance de l'air, de la rotation terrestre et de la température\footnote{Le calcul d'une seule trajectoire demandait environ 750 multiplications. Avec un \emph{differential analyzer} mécanique, le temps de calcul était d'environ 20 minutes par trajectoire ; à la main, une calculatrice humaine (\emph{computer}, au sens d'alors) y consacrait plusieurs jours.}. Le volume de travail dépasse la capacité de toute organisation humaine. Le laboratoire emploie alors plus d'une centaine de calculateurs, pour l'essentiel des femmes titulaires de diplômes en mathématiques, dont le travail ne suffit pas à répondre au rythme des demandes du front.

\dimmark{D8}{+}%
C'est dans ce contexte que John Mauchly et J.~Presper Eckert, à la Moore School of Electrical Engineering de l'université de Pennsylvanie, proposent en 1942 une machine entièrement électronique capable d'exécuter les mêmes calculs. L'armée finance le projet en avril 1943, pour un coût final de 400\,000 dollars, soit l'équivalent d'environ 7 millions de dollars de 2024. La machine, achevée en février 1946, arrive trop tard pour la guerre qu'elle devait servir. Le point analytique est ailleurs : seul un acteur non marchand, confronté à un besoin opérationnel sans substitut et disposant de ressources sans contrainte de rentabilité, peut financer un prototype dont le coût unitaire exclut tout marché civil. Le pattern est identique à celui que l'on a observé pour le télégraphe, lorsque le Congrès américain vote en 1843 les 30\,000 dollars de la ligne Washington-Baltimore, et pour la tabulatrice, lorsque le Census Bureau finance le concours de 1888 qui donne naissance à la machine de Hollerith. L'État crée le marché avant que le marché n'existe.

\dimmark{D5}{*}%
\textcite{flamm_creating_1988} documente cette structure avec une précision quantitative que la littérature d'histoire des techniques n'atteint pas toujours. Dans les années 1950, le gouvernement fédéral finance environ les deux tiers de la recherche et du développement en informatique aux États-Unis. Lorsque IBM décide de lancer la production en série de l'IBM 650 au milieu de la décennie, la vente projetée de cinquante machines au gouvernement fédéral, sur un total prévisionnel de deux cent cinquante unités, pèse de façon déterminante dans la décision. La demande militaire ne se limite pas au financement direct : la Guerre froide, l'explosion de la première bombe atomique soviétique en 1949, puis la guerre de Corée accélèrent les achats fédéraux de matériel informatique et poussent l'administration Truman à investir massivement dans la technologie de calcul. L'Atomic Energy Commission, la National Security Agency et le Department of Defense soutiennent chacun le développement de systèmes de calcul avancés pour la conception d'armes nucléaires, la cryptanalyse et la défense aérienne.

Le recul de la part publique dans le financement est lui-même une mesure de la transition vers le marché civil. \textcite{flamm_targeting_1987} montre que cette part passe d'environ deux tiers dans les années 1950 à moins d'un cinquième dans les années 1980, à mesure que les ventes privées d'IBM et de ses concurrents prennent le relais. La trajectoire n'est pas linéaire ; elle est marquée par des paliers qui correspondent aux grandes commandes fédérales (le projet SAGE de défense aérienne au milieu des années 1950, les programmes spatiaux de la NASA dans les années 1960). Chaque palier de financement public crée un seuil de capacité industrielle que le marché civil exploite ensuite, sans jamais avoir eu à supporter le risque initial.

\dimmark{0-TER}{+}%
Le projet SAGE (\emph{Semi-Automatic Ground Environment}, 1952--1963) est le cas le plus éloquent de ce couplage militaire entre informatique et énergie. Destiné à détecter et intercepter les bombardiers soviétiques via un réseau de vingt-trois centres reliés à des stations radar à travers l'Amérique du Nord, chaque centre SAGE est équipé d'un ordinateur AN/FSQ-7 de deux cent cinquante tonnes qui consomme trois mégawatts~--- vingt fois la consommation de l'ENIAC\footnote{Les données de puissance sont documentées dans les sources primaires~: \og~It weighed half a million pounds and consumed three megawatts of power\fg{} (interactions.acm.org, 2010, \og Project SAGE, a half-century on\fg{}). Deux machines AN/FSQ-7 étaient construites par centre, l'une en fonctionnement et l'autre en veille chaude permanente~--- une contrainte de disponibilité continue qui fixe la consommation à six mégawatts par centre, indépendamment de la charge de calcul.}. L'opération exige des générateurs spécialisés, des systèmes de refroidissement dédiés, et des bâtiments conçus autour des contraintes électriques de la machine. Western Electric est retenu comme contractant principal pour les bâtiments et l'alimentation électrique interne~--- au même niveau qu'IBM pour le hardware et la RAND Corporation pour le software. Pour la première fois dans l'histoire de l'informatique, l'ingénierie électrique n'est pas un service externe~: elle est une composante d'architecture à part entière, définie simultanément avec la machine et non après elle. Le couplage ICT/énergie cesse d'être un problème opérationnel pour devenir un problème de conception.

\dimmark{D8/D10}{+}%
Deux moments conceptuels prolongent SAGE et préparent la décennie suivante. En avril 1963, J.~C.~R. Licklider, alors directeur du \emph{Information Processing Techniques Office} de l'ARPA, distribue à ses collaborateurs un mémorandum dans lequel il propose le développement d'un \emph{intergalactic computer network} permettant l'opération intégrée de ressources de calcul géographiquement réparties~\parencite{licklider_memorandum_1963}. La même année, Paul Baran publie pour la RAND Corporation un rapport décrivant un réseau de communication distribué fondé sur la commutation par paquets, conçu pour rester opérationnel après une frappe nucléaire qui détruirait les nœuds centraux d'un réseau classique~\parencite{baran_distributed_1963}. Les deux propositions partagent un contexte~--- la défense Cold War et la conviction que la fragilité du réseau commuté par circuit imposait une architecture nouvelle~--- mais leur articulation effective demandait des moyens institutionnels dont aucun des auteurs ne disposait à ce moment. Ce que l'ARPA construira à partir de 1965~--- l'ARPANET, opérationnel en 1969 entre quatre nœuds universitaires~--- combinera la vision distribuée de Baran et l'ambition d'intégration de Licklider~\parencite{abbate_inventing_1999}. Cette généalogie est analytiquement importante pour la suite de la thèse~: le réseau dont naîtra Internet n'est pas un héritage civil contraint par la défense, c'est une commande de la défense dont la séquence étatique de mise en œuvre suit la même logique que celle de SAGE, du télégraphe Morse de 1843 et de la tabulatrice Hollerith de 1888. \textcite{edwards_closed_1996} a proposé de cette période une lecture canonique~: les ordinateurs de la Cold War ne sont pas seulement des artefacts techniques mais des dispositifs discursifs qui structurent la pensée du contrôle, du commandement et de l'ennemi~--- ce que l'analyse présente prolonge en y ajoutant le compte des flux matériels et financiers que la lecture discursive d'Edwards laisse hors champ.

\emergence{Le pattern demande militaire $\rightarrow$ prototype financé par l'État $\rightarrow$ marché civil n'est pas propre au computing. Il reproduit la séquence observée pour le télégraphe (§\,\ref{sec:telegraph}) et pour la tabulatrice (§\ref{subsec:hollerith}), mais à une échelle de financement incomparable. Flamm documente que le gouvernement fédéral finance les deux tiers de la R\&D informatique dans les années 1950, contre un simple vote de crédits ponctuels pour Morse en 1843. L'ampleur du financement reflète l'ampleur de la demande : le calcul balistique, la cryptanalyse et la simulation nucléaire posent des problèmes de taille $n$ tels qu'aucune organisation de travail humain, si ingénieuse soit-elle, ne peut les résoudre dans les délais opérationnels. La machine n'est pas ici un substitut au travail qualifié, comme le calculateur Burroughs remplaçant le comptable ; elle est la condition de possibilité d'un calcul qui n'existait pas avant elle.}

\questionch{2}{Comment la transition de la demande militaire à la demande civile modifie-t-elle la structure des coûts et la trajectoire technologique du computing~? Le chapitre~2 devra formaliser le mécanisme par lequel le financement public sans contrainte de rentabilité crée un seuil technologique que le marché civil ne pourrait pas atteindre seul (D8), et par lequel la demande dérivée de la défense oriente la trajectoire de l'innovation dans des directions qui ne sont pas celles que le marché aurait choisies (D5, D10).}


% -----------------------------------------------------------------------------
\subsection{Les architectures alternatives bloquées (1959--1973)}\label{subsec:alternatives_bloquees}
% -----------------------------------------------------------------------------

\dimmark{D8/D10}{+}%
Le récit qui précède pourrait laisser entendre que la trajectoire ARPANET~--- décentralisation par commutation par paquets, financement militaire américain, ouverture progressive aux universités puis au marché~--- était la seule architecture computationnelle viable au tournant des années 1960. Deux cas contemporains et un troisième, moins connu, montrent qu'il n'en était rien. Le présent chapitre les traite avec la même attention que les cas d'émergence réussie, parce qu'ils révèlent les conditions nécessaires par leur absence, en symétrie avec ce qui a été établi au XIX\textsuperscript{e}~siècle pour la \emph{Difference Engine} de Babbage et pour le télégraphe électrique français entre 1837 et 1851 (§\ref{sec:telegraph}).

\dimmark{D8}{+}%
En Union soviétique, le mathématicien Viktor Glouchkov propose à partir de 1959 le projet OGAS (\emph{Общегосударственная автоматизированная система}, Système national automatisé), conçu pour articuler par un réseau d'ordinateurs les flux d'information de planification entre les ministères, les régions et les usines de l'économie planifiée. \textcite{peters_how_2016} documente le trajet du projet sur trois décennies (1959--1989) et son échec final, qu'il n'attribue pas à des limitations techniques~--- l'URSS dispose alors d'une informatique générale compétente, dont l'ordinateur Oural (1955--1990) documente la diffusion industrielle dans une lecture complémentaire~\parencite{koretsky_ural_2022}~--- mais à la concurrence bureaucratique entre ministères, dont chacun craignait la perte de contrôle informationnel qu'aurait imposé un système intégré. Le titre choisi par Peters pour son ouvrage~--- \emph{How Not to Network a Nation}~--- circonscrit l'objet de la démonstration~: ce n'est pas que la nation n'a pas pu être mise en réseau, c'est que les conditions politiques et institutionnelles d'une telle mise en réseau n'étaient pas réunies, et l'explication des conditions absentes est aussi informative que celle des conditions réunies aux États-Unis pour ARPANET.

\dimmark{D8/0-GOV}{+}%
Au Chili, le gouvernement Allende lance en 1971 le projet Cybersyn (\emph{Synco}), conçu par le cybernéticien britannique Stafford Beer et soutenu directement par Fernando Flores au ministère de l'Économie~: un réseau de Telex reliant les usines nationalisées à un centre de coordination équipé d'un IBM 360/50, permettant le pilotage de l'économie socialiste en temps quasi réel par un tableau de bord cybernétique. \textcite{medina_cybernetic_2011} a reconstitué l'architecture du projet, sa trajectoire courte et son contexte intellectuel. Le projet est interrompu par le coup d'État du 11~septembre 1973~: la salle de contrôle est démantelée, les archives dispersées, les ingénieurs exfiltrés ou tués. La frontière entre non-émergence et destruction politique est ici prise au mot~: Cybersyn n'a pas échoué par défaut de demande ni par manque de financement, il a été supprimé par une intervention politique externe. La différence avec OGAS, dont l'échec est lent et endogène, mérite d'être préservée analytiquement~: les mécanismes de non-émergence ne sont pas tous du même type.

\dimmark{D10}{+}%
Ces cas ne sont pas des contre-factuels~: ils ne disent pas ce qui se serait passé si OGAS ou Cybersyn avaient été déployés, ni ne prétendent que ces architectures auraient été socialement préférables à celle qui a effectivement émergé. Ils disent que la trajectoire ARPANET~--- décentralisée, militaire, américaine~--- n'épuise pas l'espace des configurations computationnelles disponibles au tournant des années 1970, et que d'autres architectures, fondées sur d'autres rapports entre l'État et le calcul, ont été techniquement portées puis institutionnellement bloquées. \textcite{johnstone_deep_2026} font de cette observation un argument structurel du déploiement de la digitalisation~: la trajectoire dominante s'impose moins par sa supériorité intrinsèque que par l'absence des conditions politiques qui auraient permis l'émergence d'alternatives, et le mécanisme se reproduira lors de bifurcations ultérieures que la section~\ref{sec:convergence} examinera (refus d'AT\&T en 1971 de racheter l'ARPANET~; rejet en 1994 du projet de loi Inouye réservant vingt pour cent des capacités télécoms américaines à des usages publics). La symétrie avec la non-émergence du télégraphe électrique français entre 1837 et 1851 et avec la \emph{Difference Engine} de Babbage est analytique~: dans les trois époques, c'est la dimension institutionnelle (\textbf{0-GOV}) et la dimension financement (D8) qui déterminent la réalisation effective ou bloquée des architectures techniquement disponibles, non la supériorité intrinsèque d'une option sur les autres.

\questionch{2}{Quelles conditions institutionnelles (\textbf{0-GOV}) et financières (D8) doivent être réunies pour qu'une architecture computationnelle techniquement portée atteigne effectivement le déploiement~? Et symétriquement~: quels mécanismes bloquent une architecture alternative quand même ses conditions techniques sont réunies~? Le chapitre~2 devra formaliser ces mécanismes en distinguant les non-émergences endogènes (OGAS, fragmentation bureaucratique) des non-émergences par destruction externe (Cybersyn, coup d'État).}


% -----------------------------------------------------------------------------
\subsection{La loi de Moore et l'économie politique du semi-conducteur}\label{subsec:moore_law}
% -----------------------------------------------------------------------------

\dimmark{D10}{*}%
En~1965, Gordon Moore, cofondateur d'Intel, publie dans \emph{Electronics} un article de quatre pages qui deviendra le texte le plus cité de l'histoire de l'industrie électronique \parencite{moore_cramming_1965}. L'observation est empirique~: le nombre de composants par circuit intégré, à coût unitaire minimal, double approximativement tous les dix-huit mois. Mais la force du texte ne réside pas dans l'observation~; elle réside dans sa fonction performative. Moore lui-même l'a formulé sans ambiguïté~: \og \emph{My prediction was about the future direction of the semiconductor industry, and I have found that the industry is best understood through some of its underlying economics}\fg{}\footnote{Cité par \textcite{roussilhe_phase_2025}, qui s'appuie sur Babbage (2024), \og Moore on Moore\fg{}, \emph{The Chip Letter}. \textcite{mody2017moores} documente en détail la construction sociale de la \og loi\fg{} et son rôle dans l'organisation de la R\&D du secteur.}. La \og loi\fg{} de Moore n'est pas une loi physique~; c'est un dispositif de coordination sectorielle qui aligne les anticipations des fabricants de puces, des éditeurs de logiciel et des investisseurs sur une feuille de route commune. Tant que l'industrie croit en sa capacité à maintenir le rythme, elle investit les sommes nécessaires pour le maintenir~; si elle cessait d'y croire, les investissements tariraient et la \og loi\fg{} s'arrêterait. Comme le formule \textcite{roussilhe_phase_2025}, c'est un \og baromètre de la croyance de l'industrie en elle-même\fg{}.

\dimmark{D8}{+}%
Cette croyance a un coût. La \og seconde loi de Moore\fg{}, ou loi de Rock, stipule que le prix d'une fonderie de semi-conducteurs de nouvelle génération double tous les quatre ans \parencite{ross_moores_1995}. Le coût de construction d'une fonderie capable de graver en dessous de 10~nanomètres dépasse désormais vingt milliards de dollars, ce qui exclut de fait toutes les entreprises sauf trois ou quatre dans le monde \parencite{gruber_semiconductor_1996, roussilhe_phase_2025}. L'écart entre les deux \og lois\fg{} définit la condition de survie de l'industrie~: il faut que les économies d'échelle obtenues par la miniaturisation (première loi) croissent plus vite que les coûts de production (seconde loi). Quand cet écart se resserre, la pression à \og consommer\fg{} la puissance de calcul mise sur le marché s'intensifie~--- un mécanisme dont les conséquences sont traitées dans la section~\ref{sec:convergence}.

\dimmark{0-TER}{*}%
L'évolution de la puissance électrique maximale des composants (TDP, \emph{Thermal Design Power}) fournit un proxy factuel de la trajectoire énergétique. L'Intel 4004, premier microprocesseur commercialisé en~1971, affichait un TDP de 0,63~watt. Le Motorola~68000 (1979) atteignait 1,5~watt~; l'Intel Pentium (1994), 11,7~watts~; le Pentium~4 Extreme Edition (2003), 92~watts \parencite{sun_chip_2024, roussilhe_phase_2025}. Entre~2005 et~2018, le TDP des processeurs grand public se stabilisa entre 75 et 150~watts~--- la période où le \emph{scaling} de Dennard, décrit dans la section précédente (§\,\ref{sec:electrification}), contenait encore la consommation. La suite de cette série~--- l'explosion des TDP avec les GPU dédiés à l'IA générative~--- relève de la section~\ref{sec:contemporain}.


% -----------------------------------------------------------------------------
\subsection{Le \emph{consent decree} et la naissance du logiciel comme marchandise}\label{subsec:consent_decree}
% -----------------------------------------------------------------------------

\dimmark{D2/0-GOV}{*}%
Le 14~janvier 1949, le Department of Justice dépose une plainte antitrust contre AT\&T, accusant le monopole téléphonique d'avoir utilisé sa position dominante pour exclure les concurrents du marché des équipements. L'affaire se conclut le 24~janvier 1956 par un \emph{consent decree} dont les conséquences dépassent largement le secteur des télécommunications. AT\&T conserve son monopole sur le service téléphonique, mais doit licencier l'ensemble de ses brevets~--- y compris le transistor~--- à tout demandeur, moyennant des redevances raisonnables. Texas Instruments, Fairchild Semiconductor et la future Intel bénéficient directement de cette diffusion forcée~: la technologie du transistor, développée aux frais des abonnés téléphoniques, devient accessible à l'industrie électronique tout entière\footnote{\textcite{temin_fall_1987} documente l'histoire du \emph{consent decree} de 1956 et ses effets sur la structure de l'industrie américaine des télécommunications. La clause de licence obligatoire est l'un des mécanismes par lesquels l'antitrust américain a involontairement accéléré la diffusion des technologies de l'information.}.

\dimmark{D2/D3}{*}%
Treize ans plus tard, le mécanisme se reproduit pour IBM. En janvier 1969, le dernier jour de l'administration Johnson, le Department of Justice ouvre une enquête antitrust contre IBM, accusant la firme de monopoliser le marché des ordinateurs en liant (\emph{bundling}) matériel, logiciel et services dans une offre indissociable. Le 23~juin 1969, six mois après l'ouverture de l'enquête et avant même toute décision judiciaire, IBM annonce l'\emph{unbundling}~: désormais, le matériel, le logiciel et les services seront facturés séparément. La décision est présentée comme volontaire~; elle anticipe une contrainte judiciaire que la firme préfère devancer. L'affaire elle-même traînera jusqu'en 1982, date à laquelle l'administration Reagan l'abandonnera~; mais le mal est fait, du point de vue d'IBM. Le logiciel est devenu une marchandise autonome.

\dimmark{D3}{*}%
Cette séparation a une conséquence comptable qui prépare la section~\ref{sec:convergence}. Avant 1969, le logiciel n'existe pas comme catégorie économique distincte~: il est inclus dans le prix de la machine, sans valeur comptable propre. Après 1969, le logiciel devient un bien séparé, appropriable, valorisable. Les premières entreprises de logiciel indépendant~--- Computer Associates (1976), SAP (1972), Oracle (1977)~--- naissent dans le sillage de l'\emph{unbundling}. \textcite{campbell-kelly_software_2003} documente cette naissance comme la création d'une industrie qui n'existait pas dix ans plus tôt. Le capital intangible de \textcite{corrado_intangible_2005}, dont la mesure structure la comptabilité nationale depuis les années 2000, trouve son origine institutionnelle dans une décision antitrust de 1969.

\dimmark{D3/0-TER}{+}%
Le point analytique pour la thèse est le suivant. L'\emph{unbundling} crée une asymétrie entre la visibilité comptable du logiciel et l'invisibilité de son substrat énergétique. Le logiciel, une fois séparé du matériel, peut être valorisé comme actif immatériel~: il figure dans les bilans, il est amorti, il génère des revenus de licence. Mais le logiciel n'existe que s'il est exécuté, et l'exécution consomme de l'électricité. Cette consommation n'est pas attachée au logiciel dans les comptes~: elle est noyée dans les charges d'exploitation de l'utilisateur, pas dans le coût de production du vendeur. L'\emph{unbundling} de 1969 inaugure ainsi une dissociation comptable entre la valeur du logiciel (visible, appropriable) et le coût énergétique de son usage (diffus, non imputable), dissociation qui structure encore aujourd'hui la comptabilité du secteur numérique.


% -----------------------------------------------------------------------------
\subsection{Bilan~: trois ruptures et un paradoxe}\label{subsec:bilan_computing}
% -----------------------------------------------------------------------------

\dimmark{0-TER}{*}%
La période 1945--1975 voit se produire trois ruptures simultanées dont la conjonction prépare les développements des sections suivantes.

\paragraph{Première rupture~: l'architecture à programme enregistré.} Le \emph{stored-program concept}, formalisé par von Neumann en 1945 et matérialisé dans l'EDVAC puis l'IAS Machine, unifie l'espace mémoire du programme et l'espace mémoire des données. Le programme devient donnée~: il peut être copié, modifié, transmis comme n'importe quel fichier. Cette fusion est la condition technique de l'industrie du logiciel~: sans elle, le logiciel resterait câblé dans la machine et ne pourrait pas être vendu séparément. Mais elle a aussi une conséquence énergétique~: la mémoire qui contient le programme consomme de l'énergie en permanence (rafraîchissement DRAM). Stocker un programme n'est plus gratuit énergétiquement, contrairement au carton perforé ou au câblage fixe.

\paragraph{Deuxième rupture~: la transition tube-transistor-IC.} L'efficacité énergétique par opération s'améliore de neuf ordres de grandeur entre l'ENIAC (1946) et l'Intel 4004 (1971). Cette amélioration résulte de la miniaturisation (moins de matière à commuter, moins d'énergie dissipée) et du passage à l'électronique à semi-conducteurs (pas de filament à chauffer). Le tableau~\ref{tab:efficacite-calcul} documente les jalons de cette trajectoire.

\begin{table}[htbp]
\caption{Efficacité énergétique du calcul~: jalons 1946--1971.}\label{tab:efficacite-calcul}
\centering\small
\renewcommand{\arraystretch}{1.15}
\begin{tabular}{@{}llrrr@{}}
\toprule
\textbf{Année} & \textbf{Machine} & \textbf{Puissance (kW)} & \textbf{Opérations/s} & \textbf{Ops/kWh} \\
\midrule
1946 & ENIAC & 150 & 5\,000 & $1{,}2 \times 10^{2}$ \\
1954 & IBM 704 & 30 & 40\,000 & $4{,}8 \times 10^{3}$ \\
1964 & IBM 360/50 & 20 & 500\,000 & $9{,}0 \times 10^{4}$ \\
1971 & Intel 4004 & 0,0005 & 60\,000 & $4{,}3 \times 10^{11}$ \\
\bottomrule
\end{tabular}
\renewcommand{\arraystretch}{1.0}
\par\smallskip
\footnotesize Sources~: \textcite{nordhaus_two_2007}, \textcite{koomey_implications_2011}. Ops/kWh = opérations par kilowattheure.
\end{table}

\paragraph{Troisième rupture~: l'\emph{unbundling} et la naissance du capital intangible.} L'\emph{unbundling} d'IBM (1969) crée le logiciel comme marchandise séparée du matériel. Cette séparation institutionnelle prépare la mesure du capital intangible de \textcite{corrado_intangible_2005}, qui structure aujourd'hui la comptabilité nationale. Elle prépare aussi la dissociation comptable entre la valeur du logiciel (visible, appropriable) et le coût énergétique de son usage (diffus, non imputable).

\dimmark{0-TER}{*}%
\paragraph{Le paradoxe.} Ces trois ruptures ont un effet combiné paradoxal. L'efficacité par opération s'améliore spectaculairement~; la consommation agrégée croît parce que le volume de calcul croît plus vite que l'efficacité. Le même mécanisme de rebond (\emph{rebound effect}) que \textcite{jevons_coal_1865} avait identifié pour le charbon s'applique au calcul~: chaque gain d'efficacité libère des ressources qui sont réinvesties en volume, de sorte que la consommation totale augmente. La section~\ref{sec:convergence} documentera comment ce paradoxe se déploie entre 1975 et 2007~; la section~\ref{sec:contemporain} montrera qu'il s'intensifie après 2007 avec l'IA générative.

\emergence{La période 1945--1975 voit trois ruptures simultanées~: l'architecture à programme enregistré qui fusionne STORE et COMPUTE~; la transition tube-transistor-IC qui améliore l'efficacité énergétique de neuf ordres de grandeur~; l'\emph{unbundling} qui crée le logiciel comme capital intangible séparé du matériel. Ces ruptures préparent la convergence numérique de §\,\ref{sec:convergence} et l'explosion de §\,\ref{sec:contemporain}. Le paradoxe qui les traverse~--- amélioration de l'efficacité unitaire, croissance de la consommation agrégée~--- est le mécanisme central de l'occultation de la dimension 0-TER dans les décennies qui suivent.}

\questionch{3}{Comment quantifier le rebond énergétique du calcul sur la période 1945--2020~? Le croisement des séries d'efficacité de \textcite{koomey_implications_2011} avec les séries de capacité de \textcite{hilbert_worlds_2011} et les statistiques de consommation électrique des datacenters de \textcite{masanet_recalibrating_2020} permettrait de tester l'hypothèse Jevons sur le calcul. Ce croisement n'a jamais été publié sous cette forme.}


\begin{table}[htbp]
\caption{Dimensions économiques du calcul programmable (§\ref{sec:computing}).}\label{tab:computing-dimensions}
\centering\small
\renewcommand{\arraystretch}{1.15}
\begin{tabularx}{\textwidth}{@{}cXc@{}}
\toprule
\textbf{Dim.} & \textbf{Ce que §\ref{sec:computing} établit} & \textbf{Niveau} \\
\midrule
D2 & \emph{Consent decree} 1956 (AT\&T) et \emph{unbundling} 1969 (IBM)~: l'antitrust comme mécanisme de diffusion technologique. & ★ \\
D3 & \emph{Unbundling} 1969~: naissance du logiciel comme capital intangible séparé du matériel. & ★ \\
D5 & Demande militaire (tables de tir, cryptanalyse, simulation nucléaire, SAGE)~: l'État crée le marché que le privé ne pourrait financer. & ★ \\
D8 & L'État finance 2/3 de la R\&D dans les années 1950~; le marché civil prend le relais dans les années 1970--1980. & ★ \\
D10 & Loi de Moore comme dispositif de coordination performatif~; \emph{scaling} de Dennard comme bouclier d'occultation énergétique. Architectures alternatives bloquées (OGAS, Cybersyn). & ★ \\
\midrule
0-TER & Efficacité $\times 10^{9}$, mais consommation agrégée croît (rebond Jevons). Architecture von Neumann couple STORE et énergie (rafraîchissement DRAM). ENIAC = 150\,kW~; SAGE = 3\,MW par centre. & ★ \\
0-GOV & Antitrust (consent decree, unbundling)~; financement fédéral massif~; trajectoires alternatives bloquées par la configuration politique (OGAS, Cybersyn). & ★ \\
\bottomrule
\end{tabularx}
\renewcommand{\arraystretch}{1.0}
\end{table}